% Part: lambda-calculus
% Chapter: introduction
% Section: composition

\documentclass[../../../include/open-logic-section]{subfiles}

\begin{document}

\olfileid{lam}{int}{com}
\olsection{\usetoken{S}{lambda definable} ప్రమేయాలు సంయుక్తం కింద సంవృతమైనవి}

\begin{lem}
\tetoken{లాంబ్డాతో నిర్వచించదగిన}{!!{lambda definable}} ప్రమేయాలు సంయుక్తం కింద సంవృతమైనవి.
\end{lem}

\begin{proof}
$f$ను $h$, $g_0,$ \dots,~$g_{k-1}$ల నుంచి సంయుక్తం ద్వారా
నిర్వచించారని అనుకుందాం. $h$, $g_0$, \dots,~$g_{k-1}$లు వరుసగా
$H$, $G_0$, \dots,~$G_{k-1}$లచే \tetoken{లాంబ్డాతో నిర్వచించబడిన}{!!{lambda defined}}వని అనుకుంటే,
పదం~$F$ను కనుగొనాలి; అది $f$ను \tetoken{లాంబ్డాతో నిర్వచిస్తుంది}{!!{lambda define}s}. కానీ $F$ను కేవలం
\[
F(x_0, \dots, x_{l-1}) = H(G_0(x_0, \dots, x_{l-1}),
\dots, G_{k-1}(x_0, \dots, x_{l-1})).
\]
ద్వారా నిర్వచించవచ్చు. మరో మాటలో, లాంబ్డా కలనశాస్త్ర భాష సంయుక్తానికి
ప్రాతినిధ్యం వహించడానికి చక్కగా సరిపోతుంది.
\end{proof}

\end{document}
